\section{Mekanisme Deteksi Drift}

Pergeseran distribusi antar-jendela diukur sebagai Jensen-Shannon divergence atas marginal pasangan atribut-nilai; struktur dependensi hanya dipelajari ulang ketika pergeseran melampaui ambang $\tau$, sedangkan parameter (\textit{Probability Vector}) tetap diperbarui tiap siklus.

\begin{verbatim}
# distribusi frekuensi pasangan atribut-nilai 
# pada satu jendela
def av_marginal(A):     
    m = A.mean(axis=0).astype(np.float64); s = m.sum()
    return m / s if s > 0 else np.full(len(m), 
    1.0 / max(1, len(m)))

# Jensen-Shannon divergence
def jsd(p, q, eps=1e-12):                       
    m  = 0.5 * (p + q)
    kl = lambda a: float(np.sum(np.where(a > 0, 
    a * np.log2((a+eps)/(m+eps)), 0.0)))
    return max(0.0, min(1.0, 
    0.5 * kl(p) + 0.5 * kl(q)))

# Pemicu (di run_sequential): struktur dipelajari 
# ulang HANYA bila drift > tau

# distribusi jendela berjalan
cur     = av_marginal(prob.A)

# pergeseran thd jendela referensi
delta   = 0.0 if ref is None else jsd(ref, cur)

# tau = 0,01 (baku)
relearn = (mode == "always") or (delta > tau)   

if relearn:
    # perbarui struktur; jika tidak: PV saja
    T = chow_liu_prim(prob.A, rng)              
\end{verbatim}

